\chapter{Session（会话）}

\vspace{0.5em}
\noindent\textit{上一章我们了解了 Binding 如何将消息路由到正确的 Agent。本章将讲解 Session（会话）——
用户与智能体之间一次连续对话的容器。我们将了解会话的生命周期、分组策略、
重置机制以及跨渠道身份统一。}
\vspace{0.5em}

\section{概述}

Session（会话）是用户与智能体之间一次连续对话的容器，包含上下文历史、令牌计数和路由元数据。

配置路径：\texttt{session.*}

\section{会话所有权}

所有会话状态都\textbf{由 Gateway 网关拥有}。UI 客户端（macOS 应用、WebChat 等）必须向 Gateway 网关查询，而不是读取本地文件。

\subsection{存储位置}

\begin{itemize}
  \item 存储文件：\texttt{\textasciitilde{}/.openclaw/agents/<id>/sessions/sessions.json}
  \item 对话记录：\texttt{\textasciitilde{}/.openclaw/agents/<id>/sessions/<sessionId>.jsonl}
\end{itemize}

\section{会话键映射}

\begin{center}
\begin{tabular}{ll}
\toprule
\textbf{来源} & \textbf{会话键格式} \\
\midrule
私信（默认） & \texttt{agent:<agentId>:<mainKey>} \\
群组聊天 & \texttt{agent:<agentId>:<channel>:group:<groupId>} \\
线程/话题 & 基础键 + \texttt{:thread:<threadId>} \\
定时任务 & \texttt{cron:<jobId>} \\
Webhook & \texttt{hook:<hookId>} \\
\bottomrule
\end{tabular}
\end{center}

\section{私信分组策略（dmScope）}

\texttt{session.dmScope} 控制私信如何分组：

\begin{center}
\begin{tabular}{ll}
\toprule
\textbf{值} & \textbf{行为} \\
\midrule
\texttt{main}（默认） & 所有私信共享主会话 \\
\texttt{per-peer} & 按发送者 ID 隔离 \\
\texttt{per-channel-peer} & 按渠道 + 发送者隔离 \\
\texttt{per-account-channel-peer} & 按账户 + 渠道 + 发送者隔离 \\
\bottomrule
\end{tabular}
\end{center}

\section{会话生命周期}

\subsection{重置策略}

\begin{itemize}
  \item \textbf{每日重置}（默认）：Gateway 主机本地时间凌晨 4:00
  \item \textbf{空闲重置}（可选）：滑动空闲窗口（\texttt{idleMinutes}）
  \item 两者同时配置时，先过期者触发新会话
\end{itemize}

\subsection{手动重置}

在聊天中发送 \texttt{/new} 或 \texttt{/reset} 创建新会话。

\section{配置示例}

\begin{lstlisting}[language=json]
{
  session: {
    dmScope: "main",
    reset: {
      mode: "daily",
      atHour: 4,
      idleMinutes: 120
    },
    resetByType: {
      dm: { mode: "idle", idleMinutes: 240 },
      group: { mode: "idle", idleMinutes: 120 }
    },
    identityLinks: {
      alice: ["telegram:123456789", "discord:987654321012345678"]
    }
  }
}
\end{lstlisting}

\texttt{identityLinks} 可将同一用户的不同渠道 ID 映射到统一身份，实现跨渠道共享会话。

\section{常用命令}

\begin{lstlisting}[language=bash]
openclaw status                      # Show sessions
openclaw sessions --json             # Export all sessions
/status                              # In chat: check session info
/new                                 # In chat: start new session
/compact                             # In chat: compact context
/stop                                # In chat: abort current run
\end{lstlisting}

\section{本章小结}

本章介绍了 Session 会话管理的核心机制：

\begin{itemize}
  \item Session 是用户与智能体之间\textbf{一次连续对话的容器}，由 Gateway 统一管理
  \item \textbf{dmScope} 控制私信分组策略，从共享主会话到按发送者隔离
  \item 支持\textbf{每日重置}和\textbf{空闲重置}两种自动重置策略
  \item \texttt{identityLinks} 实现跨渠道身份统一
\end{itemize}

\noindent 到目前为止，我们已经掌握了 OpenClaw 的基础架构（Gateway、Agent、Channel、Binding、Session）。
接下来我们将进入\textbf{扩展能力}的篇章——下一章讲解 Plugin（插件），看看如何为 OpenClaw 添加新能力。
